\section{System Modeling and Problem Formulation}
\label{sec:modeling}

\subsection{Rehabilitation Scenario and Coordinate Definition}
\label{subsec:scenario}

% Trace: III-A-P1
We consider prescribed passive hip--knee flexion for a supine subject.  The
lower limb is approximated in the sagittal plane, with the hip center fixed at
the origin of a rehabilitation frame.  The $x$-axis lies along the bed and
points from the hip toward the foot, while the $z$-axis points upward.  This
model describes the thigh, knee, shank, and the point at which the robot acts
through a lower-leg strap; it does not model active muscle recruitment,
distributed strap pressure, out-of-plane motion, or robot-link dynamics.
% TODO: citation for the sagittal-plane and passive-rehabilitation assumptions.

% Trace: III-A-P2
The generalized coordinates are
\begin{equation}
    \vect{q}=\begin{bmatrix}q_h & q_k\end{bmatrix}^{\mathsf T},
    \label{eq:generalized_coordinates}
\end{equation}
where $q_h$ and $q_k$ denote hip and knee flexion, respectively.  Positive knee
flexion folds the shank toward the thigh; consequently, the absolute shank
orientation used throughout the implementation is
\begin{equation}
    \theta_s=q_h-q_k,
    \label{eq:shank_angle}
\end{equation}
not $q_h+q_k$.  All internal angular quantities are expressed in radians.

% Trace: III-A-P3
Let $L_1$ be the thigh length and let $L_2$ be the distance from the knee to the
strap equivalent traction point.  The latter is a task-specific kinematic
length and is not assumed to be the anatomical knee--ankle distance or the
shank center-of-mass distance.  The current software geometry uses
$L_1=\SI{0.42}{m}$ and $L_2=\SI{0.30}{m}$.  Hip flexion is constrained to be no
greater than $120^{\circ}$.  The repository currently uses different knee-ROM
upper limits in the legacy workspace/identification pipeline and the active
reference audit; therefore, the applicable knee interval is retained as an
explicit experiment configuration rather than asserted as a universal model
constant.

\subsection{Two-Link Lower-Limb Kinematics}
\label{subsec:kinematics}

% Trace: III-B-P1
With $\vect{p}_K=[x_K,z_K]^{\mathsf T}$ denoting the knee and
$\vect{p}_P=[x_P,z_P]^{\mathsf T}$ denoting the strap equivalent traction point,
the implemented forward kinematics are
\begin{align}
    x_K &= L_1\cos q_h, & z_K &= L_1\sin q_h, \nonumber\\
    x_P &= L_1\cos q_h+L_2\cos(q_h-q_k),
    & z_P &= L_1\sin q_h+L_2\sin(q_h-q_k).
    \label{eq:forward_kinematics}
\end{align}
The signs in~\eqref{eq:forward_kinematics} follow directly from
\eqref{eq:shank_angle}: increasing $q_k$ rotates the shank clockwise relative
to the thigh in the selected $x$--$z$ convention.

% Trace: III-B-P2
For workspace construction, a sampled posture is retained only when the knee
and traction point do not cross below the bed and the traction point does not
move behind the hip, i.e., $z_K\geq 0$, $z_P\geq 0$, and $x_P\geq 0$.  These are
geometric feasibility checks, not clinical safety guarantees.  Inverse
kinematics uses the flexion branch
\begin{align}
    D &= \frac{x_P^2+z_P^2-L_1^2-L_2^2}{2L_1L_2}, \nonumber\\
    q_k &= \cos^{-1}D, \nonumber\\
    q_h &= \operatorname{atan2}(z_P,x_P)
      +\operatorname{atan2}\!\left(L_2\sin q_k,L_1+L_2\cos q_k\right),
    \label{eq:inverse_kinematics}
\end{align}
and rejects solutions outside the configured ROM.

\subsection{Jacobian and Mechanical Interaction Mapping}
\label{subsec:jacobian}

% Trace: III-C-P1
Differentiating~\eqref{eq:forward_kinematics} gives
\begin{equation}
    \dot{\vect{p}}_P=\vect{J}(\vect{q})\dot{\vect{q}},
    \qquad
    \vect{J}(\vect{q})=
    \begin{bmatrix}
    -L_1\sin q_h-L_2\sin(q_h-q_k) & L_2\sin(q_h-q_k)\\
     L_1\cos q_h+L_2\cos(q_h-q_k) & -L_2\cos(q_h-q_k)
    \end{bmatrix}.
    \label{eq:jacobian}
\end{equation}

% Trace: III-C-P2
Under the implemented point-force assumption, let
$\vect{F}_{R\rightarrow L}=[F_x,F_z]^{\mathsf T}$ be the force exerted by the
robot on the leg at the strap equivalent traction point.  Virtual work yields
the generalized interaction torque
\begin{equation}
    \vect{\tau}_{\mathrm{meas}}
    =\vect{J}(\vect{q})^{\mathsf T}\vect{F}_{R\rightarrow L}.
    \label{eq:force_to_torque}
\end{equation}
Conversely, when a required generalized torque is mapped to an endpoint force,
the implementation evaluates
\begin{equation}
    \vect{F}_{R\rightarrow L}
    =\left(\vect{J}^{\mathsf T}\right)^{\dagger}\vect{\tau}_{\mathrm{req}},
    \label{eq:torque_to_force}
\end{equation}
where $(\cdot)^{\dagger}$ denotes the Moore--Penrose pseudoinverse.

% Trace: III-C-P3
The mapping is marked invalid when its inputs or outputs are non-finite, when
$|\det\vect{J}|<10^{-4}$, when the two-norm condition number exceeds $100$, or
when the software anomaly limit of \SI{1000}{N} is exceeded.  The last value is
only a numerical audit bound; it is not a robot or participant safety limit.
Equation~\eqref{eq:force_to_torque} also assumes that the measured force is
expressed in the modeled sagittal frame and acts at the modeled point.  Wrench
frame, sign, reference point, compensation, and physical synchronization remain
separate real-robot validation items.

\subsection{Subject-Specific Equivalent Dynamic Model}
\label{subsec:dynamics}

% Trace: III-D-P1
The modeled torque required to follow a prescribed motion is
\begin{equation}
    \vect{\tau}_{\mathrm{req}}
    =\vect{M}(\vect{q})\ddot{\vect{q}}
    +\vect{c}(\vect{q},\dot{\vect{q}})
    +\vect{g}(\vect{q})
    +\vect{B}\dot{\vect{q}}
    +\vect{K}(\vect{q}-\vect{q}_0).
    \label{eq:inverse_dynamics}
\end{equation}
Let $m_1,m_2$, $c_1,c_2$, and $I_1,I_2$ be the thigh and shank masses,
center-of-mass distances, and planar inertias.  Defining
\begin{equation}
    a=m_1c_1^2+I_1+m_2L_1^2,\quad
    b=m_2c_2^2+I_2,\quad
    d=m_2L_1c_2,
    \label{eq:dynamic_abbreviations}
\end{equation}
the implemented mass matrix is
\begin{equation}
    \vect{M}(\vect{q})=
    \begin{bmatrix}
      a+b+2d\cos q_k & -(b+d\cos q_k)\\
      -(b+d\cos q_k) & b
    \end{bmatrix}.
    \label{eq:mass_matrix}
\end{equation}
Its negative off-diagonal terms result from
$\dot{\theta}_s=\dot q_h-\dot q_k$.

% Trace: III-D-P2
The implemented Coriolis/centrifugal and gravity vectors are
\begin{align}
    \vect{c}(\vect{q},\dot{\vect{q}})
    &=d\sin q_k
    \begin{bmatrix}
      \dot q_k^2-2\dot q_h\dot q_k\\
      \dot q_h^2
    \end{bmatrix},
    \label{eq:coriolis_vector}\\
    \vect{g}(\vect{q})
    &=g
    \begin{bmatrix}
    (m_1c_1+m_2L_1)\cos q_h+m_2c_2\cos(q_h-q_k)\\
    -m_2c_2\cos(q_h-q_k)
    \end{bmatrix}.
    \label{eq:gravity_vector}
\end{align}
The negative knee component in~\eqref{eq:gravity_vector} is another consequence
of the shank-angle convention.

% Trace: III-D-P3
Passive joint behavior is represented by independent linear terms
\begin{equation}
    \vect{B}=\operatorname{diag}(b_h,b_k),\qquad
    \vect{K}=\operatorname{diag}(k_h,k_k),
    \label{eq:passive_terms}
\end{equation}
where $k_h,k_k$ are equivalent passive stiffnesses and $b_h,b_k$ are
equivalent passive damping coefficients.  The estimator changes only a common
mass/inertia scale and these four passive coefficients; segment proportions,
center-of-mass locations, neutral angles, and gravity are supplied by a
baseline template.  No passive cross-joint stiffness or damping coefficient is
estimated by the current five-parameter model.

% Trace: III-D-P4
The identified quantities are termed \emph{subject-specific equivalent
dynamics}.  They summarize the mechanical input--output relation over the
excited task region and may absorb unmodeled tissue, strap, or joint effects.
They are not interpreted as unique physiological parameters or direct tissue
measurements.  This distinction is essential when low prediction residuals are
obtained from a simplified model.

\subsection{Reference-Trajectory-Centered Personalization Problem}
\label{subsec:personalization_problem}

% Trace: III-E-P1
Let $s\in[0,1]$ be normalized cycle phase and let
$\vect{q}_{\mathrm{ref}}(s)$ be the prescribed rehabilitation reference.  The
required topological condition is
\begin{equation}
    \vect{q}_{\mathrm{ref}}(0)=\vect{q}_{\mathrm{ref}}(1).
    \label{eq:reference_closure}
\end{equation}
The path may be periodic, closed, continuous, and asymmetric; in particular,
no constraint of the form
$\vect{q}_{\mathrm{ref}}(s)=\vect{q}_{\mathrm{ref}}(1-s)$ is imposed.  The
currently active software reference retains distinct measured flexion and
extension branches and applies a small periodic cubic B-spline correction to
obtain $C^2$ closure.

% Trace: III-E-P2
A reference-local candidate can be represented abstractly as
\begin{equation}
    \vect{q}_{\mathrm{cand}}(s;\vect{\alpha})
    =\vect{q}_{\mathrm{ref}}(s)
    +\Delta\vect{q}(s;\vect{\alpha}),
    \qquad
    \Delta\vect{q}(0;\vect{\alpha})
    =\Delta\vect{q}(1;\vect{\alpha})=\vect{0}.
    \label{eq:candidate_parameterization}
\end{equation}
The repository implements fixed low-dimensional amplitude reductions, a
monotone endpoint-preserving knee-phase warp, and duration scaling.  It does
not yet implement a continuous optimizer over $\vect{\alpha}$ about the active
asymmetric reference; therefore,~\eqref{eq:candidate_parameterization} is a
problem-formulation abstraction rather than a claim of a finalized optimizer.

% Trace: III-E-P3
For a candidate set $\mathcal{Q}_{\mathrm{cand}}$, the intended task-local
selection problem is
\begin{equation}
    \vect{q}^{\star}\in
    \arg\min_{\vect{q}\in\mathcal{Q}_{\mathrm{cand}}}
    \mathcal{L}_{\mathrm{int}}(\vect{q};\hat{\vect{\vartheta}})
    \quad\text{subject to}\quad
    \vect{q}\in\mathcal{F},
    \label{eq:personalization_problem}
\end{equation}
where $\hat{\vect{\vartheta}}$ denotes the identified equivalent dynamics,
$\mathcal{L}_{\mathrm{int}}$ is a mechanical-interaction criterion, and
$\mathcal{F}$ is a feasibility set.  In the current offline candidate evaluator,
$\mathcal{F}$ checks joint ROM, $x_P\geq0$, $z_P\geq0$, $z_K\geq0$, cycle
closure, finite force mapping, Jacobian conditioning, and at least $90\%$
coverage by a train-fitted six-dimensional local state box.

% Trace: III-E-P4
Explicit velocity, acceleration, and physical interaction-load limits are
present only as planned experiment-specific constraints or in the separate
fail-closed robot preflight architecture; they are not yet validated as
personalization constraints.  Likewise, the stored candidate study reports an
unweighted multi-objective Pareto comparison, but no final subject-specific
selector, calibrated tactile criterion, or clinically validated objective.
Accordingly, the current formulation supports model-based candidate screening,
not a claim of a comfortable, clinically optimal, or clinically safer
trajectory.
